% Part: first-order-logic
% Chapter: tableaux
% Section: provability-consistency

\documentclass[../../../include/open-logic-section]{subfiles}

\begin{document}

\iftag{FOL}
      {\olfileid{fol}{tab}{prv}}
      {\olfileid{pl}{tab}{prv}}

\olsection{\usetoken{S}{derivability} మరియు అవైరుధ్యం}

ఇప్పుడు !!{derivability} సంబంధపు అనేక ధర్మాలను స్థాపిస్తాం. అవి
స్వతంత్రంగా ఆసక్తికరమైనవి; అంతేకాక ప్రతి ఒక్కటీ సంపూర్ణతా సిద్ధాంతపు
నిరూపణలో ఒక పాత్ర పోషిస్తుంది.

\begin{prop}\ollabel{prop:provability-contr}
  $\Gamma \Proves !A$ అయి, $\Gamma \cup \{!A\}$ వైరుధ్యభరితమైతే,
  $\Gamma$ వైరుధ్యభరితమైనది.
\end{prop}

\begin{proof}
$\Gamma_0 = \{!B_1, \dots, !B_n\}$,
$\Gamma_1 =\{!C_1, \dots, !C_m\} \subseteq \Gamma$ అనే పరిమిత
సమితులు ఉండి, కింది రెండింటికీ సంవృత !!{tableau}s ఉంటాయి:
\sourcecorrection{OLTETAB-005}{ఆంగ్ల మూలం $\Gamma_1$ను
$\{!C_1,\dots,!C_n\}$గా నిర్వచించి, వెంటనే అదే సమితిని
$!C_1,\dots,!C_m$గా ఉపయోగించింది. $\Gamma_0$కు $n$, $\Gamma_1$కు
$m$ అనే స్వతంత్ర పరిమిత పరిమాణాలను ఉంచేలా చివరి సూచికను $m$గా
సరిచేశాం.}
  \begin{align*}
    \{\sFmla{\False}{!A}, &
    \sFmla{\True}{!B_1}, \dots, \sFmla{\True}{!B_n}\} \\
    \{\sFmla{\True}{!A}, &
    \sFmla{\True}{!C_1}, \dots, \sFmla{\True}{!C_m}\}
  \end{align*}
  $!A$పై \Cut{} నియమాన్ని ఉపయోగించి వీటిని ఒకే సంవృత
  !!{tableau}గా కలపవచ్చు; అది $\Gamma_0 \cup \Gamma_1$
  వైరుధ్యభరితమని చూపుతుంది. $\Gamma_0 \subseteq \Gamma$,
  $\Gamma_1 \subseteq \Gamma$ కనుక $\Gamma_0 \cup
  \Gamma_1 \subseteq \Gamma$; అందువల్ల $\Gamma$ వైరుధ్యభరితమైనది.
\end{proof}

\begin{prop}
\ollabel{prop:prov-incons}
$\Gamma \Proves !A$ అయితే, అప్పుడు మాత్రమే $\Gamma \cup \{\lnot !A\}$
వైరుధ్యభరితమైనది.
\end{prop}

\begin{proof}
మొదట $\Gamma \Proves !A$ అనుకుందాం; అంటే కింది సమితికి ఒక సంవృత
!!{tableau} ఉంది:
\[
\{\sFmla{\False}{!A},
\sFmla{\True}{!B_1}, \dots, \sFmla{\True}{!B_n}\}
\]
$\TRule{\True}{\lnot}$ నియమాన్ని ఉపయోగించి దాన్ని కింది సమితికి
ఒక సంవృత !!{tableau}గా మార్చవచ్చు:
\[
\{\sFmla{\True}{\lnot !A},
\sFmla{\True}{!B_1}, \dots, \sFmla{\True}{!B_n}\}.
\]

మరోవైపు, రెండవ సమితికి ఒక సంవృత !!{tableau} ఉంటే, మొదటి పరికల్పన
$\sFmla{\True}{\lnot !A}$కు \TRule{\True}{\lnot}ను వర్తింపజేయడం
వల్ల వచ్చిన ప్రతి సూత్రాన్నీ, ఆ పరికల్పననూ తొలగించి, పరికల్పన
$\sFmla{\False}{!A}$ను చేర్చడం ద్వారా మొదటి సమితికి ఒక సంవృత
!!{tableau}గా మార్చవచ్చు. పాత శాఖ
$\sFmla{\True}{\lnot !A}$కు \TRule{\True}{\lnot}ను వర్తింపజేసిన
నిష్కర్ష, అంటే $\sFmla{\False}{!A}$ వల్ల సంవృతమైతే, కొత్త
!!{tableau}లోని సంబంధిత శాఖ కూడా సంవృతమే. పాత !!{tableau}లోని ఒక శాఖ
$\sFmla{\True}{\lnot !A}$ పరికల్పనతో పాటు
$\sFmla{\False}{\lnot !A}$ను కలిగి ఉండడం వల్ల సంవృతమైతే,
$\sFmla{\False}{\lnot !A}$కు $\TRule{\False}{\lnot}$ను
వర్తింపజేసి $\sFmla{\True}{!A}$ను పొంది దాన్ని సంవృత శాఖగా
మార్చవచ్చు. మనం $\sFmla{\False}{!A}$ను పరికల్పనగా చేర్చాం కనుక ఆ
శాఖ సంవృతమవుతుంది.
\end{proof}

\begin{prob}
$\Gamma \Proves \lnot !A$ అయితే, అప్పుడు మాత్రమే
$\Gamma \cup \{!A\}$ వైరుధ్యభరితమని నిరూపించండి.
\end{prob}

\begin{prop}\ollabel{prop:explicit-inc}
  $\Gamma \Proves !A$ అయి, $\lnot !A \in \Gamma$ అయితే $\Gamma$
  వైరుధ్యభరితమైనది.
\end{prop}

\begin{proof}
  $\Gamma \Proves !A$, $\lnot !A \in \Gamma$ అనుకుందాం. అప్పుడు
  $!B_1$, \dots, $!B_n \in \Gamma$ ఉండి, కింది సమితికి ఒక సంవృత
  !!{tableau} ఉంటుంది:
  \[
  \{\sFmla{\False}{!A}, \sFmla{\True}{!B_1}, \dots, \sFmla{\True}{!B_n}\}
  \]
  పరికల్పన \sFmla{\False}{!A} స్థానంలో
  \sFmla{\True}{\lnot !A}ను ఉంచి, అన్ని పరికల్పనల తర్వాత ఆ
  \sFmla{\True}{\lnot !A}కు \TRule{\True}{\lnot}ను వర్తింపజేసి
  వచ్చే \sFmla{\False}{!A}ను చేర్చండి. పాత !!{tableau}లో వరుస~$1$ను
  ఆశ్రయించి సమర్థించిన ఏ !!{sentence} అయినా, ఇప్పుడు కొత్తగా చేర్చిన
  ఆ వరుసను ఆశ్రయించి సమర్థించబడుతుంది. కనుక పాత !!{tableau} సంవృతమైతే
  కొత్తదీ సంవృతమే. పరికల్పనలన్నీ $\Gamma$లో ఉన్నాయి కనుక అది
  $\Gamma$ వైరుధ్యభరితమని చూపుతుంది.
  \sourcecorrection{OLTETAB-006}{ఈ మార్పు వివరణలో ఆంగ్ల మూలం
  $\TRule{\True}{\lnot}$ను తప్పుగా $\sFmla{\False}{!A}$కు
  వర్తింపజేసిందనీ, కొత్త నిష్కర్షకు పాత వరుస సూచనలు ఒక నిర్దిష్టమైన
  కానీ సరిపోని వరుస~$n+1$ను చూపాలనీ చెప్పింది. నియమ నిర్వచనం ప్రకారం
  దాన్ని $\sFmla{\True}{\lnot !A}$కు వర్తింపజేసి
  $\sFmla{\False}{!A}$ను పొందాలి; పరికల్పనల సంఖ్యపై ఆధారపడని
  “కొత్తగా చేర్చిన వరుస” సూచనతో పాత ఆధారాలను మార్చాలి. మూలంలోని
  ఒంటరి అదనపు బ్యాక్‌స్లాష్‌నూ తొలగించాం.}
\end{proof}

\begin{prop}\ollabel{prop:provability-exhaustive}
  $\Gamma \cup \{!A\}$, $\Gamma \cup \{\lnot !A\}$ రెండూ
  వైరుధ్యభరితమైతే, $\Gamma$ వైరుధ్యభరితమైనది.
\end{prop}

\begin{proof}
  $!B_1$, \dots, $!B_n \in \Gamma$, $!C_1$, \dots,
  $!C_m \in \Gamma$ ఉండి, కింది రెండు సమితులకూ సంవృత
  !!{tableau}s ఉన్నాయని అనుకుందాం:
  \begin{align*}
    \{\sFmla{\True}{!A}, &
    \sFmla{\True}{!B_1}, \dots, \sFmla{\True}{!B_n}\} \text{ and}\\
    \{\sFmla{\True}{\lnot !A}, &
    \sFmla{\True}{!C_1}, \dots, \sFmla{\True}{!C_m}\}
  \end{align*}
  $\sFmla{\True}{!B_1}$, \dots, $\sFmla{\True}{!B_n}$లతో పాటు
  $\sFmla{\True}{!C_1}$, \dots, $\sFmla{\True}{!C_m}$లను
  పరికల్పనలుగా తీసుకుని, ఆ తర్వాత \Cut{} నియమాన్ని వర్తింపజేస్తే,
  $\Gamma$ వైరుధ్యభరితమని చూపే ఒకే సంయుక్త !!{tableau}ను నిర్మించవచ్చు.
  దీనివల్ల రెండు శాఖలు వస్తాయి: ఒకటి $\sFmla{\True}{!A}$తో,
  మరొకటి $\sFmla{\False}{!A}$తో మొదలవుతుంది.

  ఎడమవైపున మొదటి !!{tableau}లో పరికల్పనల కింద ఉన్న భాగాన్ని
  చేర్చండి. ఇక్కడ మొదటి !!{tableau}లోని ప్రతి పరికల్పనా,
  $\sFmla{\True}{!A}$తో సహా, లభిస్తుంది కనుక ప్రతి నియమ వర్తింపూ
  ఇప్పటికీ సరైనదే. అందువల్ల $\sFmla{\True}{!A}$ కింద ప్రతి శాఖా
  సంవృతమవుతుంది.

  కుడివైపున రెండవ !!{tableau}లో దాని పరికల్పనల కింద ఉన్న భాగాన్ని
  చేర్చండి; అయితే $\sFmla{\True}{\lnot !A}$కు
  $\TRule{\True}{\lnot}$ను వర్తింపజేసిన అన్ని ఫలితాలను తొలగించండి.
  $\sFmla{\True}{\lnot !A}$కు $\TRule{\True}{\lnot}$ను
  వర్తింపజేస్తే వచ్చే నిష్కర్ష $\sFmla{\False}{!A}$; అది సంయుక్త
  !!{tableau}లో కుడివైపు \Cut{} నియమపు నిష్కర్షగా ఇప్పటికే లభిస్తుంది.

  రెండవ టాబ్లోలోని ఒక శాఖ ఇప్పుడు పరికల్పనగా కనిపించని
  $\sFmla{\True}{\lnot !A}$తో పాటు $\sFmla{\False}{\lnot !A}$ను
  కలిగి ఉండడం వల్ల సంవృతమైతే, $\sFmla{\False}{\lnot !A}$కు
  $\TRule{\False}{\lnot}$ను వర్తింపజేసి $\sFmla{\True}{!A}$ను
  పొందవచ్చు. సంయుక్త !!{tableau}లోని సంబంధిత శాఖ ఇప్పుడు కూడా
  సంవృతమవుతుంది; ఎందుకంటే అందులో \Cut{} నియమపు కుడివైపు నిష్కర్ష
  $\sFmla{\False}{!A}$ ఉంది. రెండవ !!{tableau}లో ఒక శాఖ మరేదైనా
  కారణంతో సంవృతమైతే, సంయుక్త !!{tableau}లోని సంబంధిత శాఖ కూడా
  సంవృతమే; పాత రెండవ !!{tableau} శాఖపై కనిపించిన
  $\sFmla{\True}{\lnot !A}$ తప్ప మిగతా ప్రతి !!{signed formula}s,
  సంయుక్త !!{tableau}లోని సంబంధిత శాఖపైనా కనిపిస్తుంది.
  \end{proof}

\end{document}
